\documentclass[11pt,a4paper]{article}

\usepackage[top=1.7cm,bottom=1.9cm,left=1.9cm,right=1.9cm]{geometry}
\usepackage{amsmath,amssymb}
\usepackage{booktabs}
\usepackage{graphicx}
\usepackage{caption}
\usepackage{subcaption}
\usepackage[hidelinks]{hyperref}
\usepackage{parskip}
\usepackage{titlesec}
\usepackage{placeins}
\usepackage{float}

\titleformat{\section}{\normalfont\large\bfseries}{\thesection}{0.6em}{}
\titleformat{\subsection}{\normalfont\normalsize\bfseries}{\thesubsection}{0.5em}{}
\titlespacing{\section}{0pt}{0.4em}{0.1em}
\titlespacing{\subsection}{0pt}{0.3em}{0.05em}

\setlength{\parskip}{2pt}
\setlength{\parindent}{0pt}

\title{\vspace{-2.5em}SF2943 Project --- Part A\\[1pt]
\large Daily electricity load for Swedish bidding zone SE\_3}
\author{}
\date{}

\begin{document}
\maketitle
\vspace{-3.6em}

\section{Data}
We use daily mean electricity load (in MW) for the Swedish bidding zone
\textbf{SE\_3} (Stockholm and the M\"alardalen region in eastern Sweden).
The data come from \emph{Open Power System Data} (OPSD), the file
\texttt{time\_series\_60min\_singleindex.csv}, column
\texttt{SE\_3\_load\_actual\_entsoe\_transparency}. The original source is the
ENTSO-E Transparency Platform.\footnote{\url{https://data.open-power-system-data.org/time_series/}}
The file contains hourly observations from 2015-01-01\,01:00 to
2020-09-30\,22:00 UTC (50\,398 hours, 77 NaN, longest gap 49~h). We aggregate
to daily means and mark days with fewer than 20 valid hours as NaN. The 5
remaining NaN days are filled by linear interpolation. This gives
$n = 2100$ daily observations $X_t$, which is well above the
$N \geq 2000$ requirement.

\section{Cleaning to stationarity}
The raw series $X_t$ has a clear winter/summer cycle, a weekday/weekend cycle,
a small long-term trend, and an amplitude that grows with the level. We use
the classical decomposition (B\&D~\S1.5) on the log of the series:
\begin{equation}
\log X_t \;=\; m_t \;+\; s_t^{(\mathrm{year})} \;+\; s_t^{(\mathrm{week})} \;+\; Y_t,
\label{eq:decomp}
\end{equation}
fitting in the order \emph{year} $\to$ \emph{week} $\to$ \emph{trend} by
ordinary least squares (OLS). Each step is fit on the residual of the previous
one. Taking logs (Box--Cox with $\lambda = 0$) stabilises the variance and
turns multiplicative seasonality into additive, which is what makes
\eqref{eq:decomp} reasonable.

\textbf{Yearly seasonality.} A $k=2$ harmonic (B\&D~\S1.3) with period $d=365$:
\[
\hat s_t^{(\mathrm{year})} \;=\; \sum_{j=1}^{2} \bigl[a_j \cos(2\pi j t / 365) + b_j \sin(2\pi j t / 365)\bigr],
\]
with $(a_1,b_1,a_2,b_2) = (+0.2384,+0.0773,-0.0127,-0.0013)$. The first
harmonic catches the main winter--summer swing and the second one handles the
asymmetric shape of the peak.

\textbf{Weekly seasonality.} A $k=3$ harmonic with period $d=7$, fit on the
year-deseasonalised series. The coefficients are
$(a_1,\dots,b_3) = (+0.0323,-0.0536,+0.0096,+0.0304,-0.0086,-0.0107)$.
Since $d=7$ is odd, $k=3$ is the saturated Fourier basis for the weekly cycle,
which is equivalent to using six day-of-week dummies.

\textbf{Polynomial trend.} We start with degree 1 on the deseasonalised series.
If the residual mean of any third of the series is bigger than 10\% of the
residual standard deviation, we move up to degree 2. Here the linear residual
thirds were $(-0.0061,+0.0113,-0.0052)$ with std $0.0645$ (max ratio
$\approx 0.18$), so we picked the quadratic
\[
\hat m_t \;=\; 9.1609 \;+\; (6.12 \times 10^{-5})\,t \;-\; (3.35 \times 10^{-8})\,t^2.
\]
The cleaning residual
$\hat Y_t = \log X_t - \hat m_t - \hat s_t^{(\mathrm{year})} - \hat s_t^{(\mathrm{week})}$
is shown in Figure~\ref{fig:decomposition}c. By eye the level and amplitude
look constant.

\textbf{Diagnostics.} Ljung--Box clearly rejects iid noise
($Q(20) = 4134$, $p \approx 0$), so cleaning has removed the deterministic
part but there is still autocorrelation left for the ARMA stage to handle.
The sample ACF (Figure~\ref{fig:diag}a) decays slowly with
$\hat\rho(1) = 0.83$, has a clear spike at lag~7 ($\hat\rho(7) = 0.30$), and
a small leftover yearly signal at $\hat\rho(365) = 0.16$. The PACF
(Figure~\ref{fig:diag}b) has its biggest spike at lag~1 with smaller activity
at lags 2--3, which suggests an AR-dominated process. As a sanity check,
$\nabla_1 \nabla_7 \log X_t$ (B\&D~\S1.5.2) gives a similar-looking ACF.

\section{ARMA fitting}
Based on the ACF/PACF shape we tried five candidate models, fit by Gaussian
MLE (B\&D~\S5.2; \texttt{statsmodels.tsa.arima.model.ARIMA} maximises the
likelihood numerically using the innovations algorithm). All five fits are
causal and invertible, i.e.\ all roots of $\hat\phi(z)$ and $\hat\theta(z)$
are outside the unit circle. We pick the order using AICC (B\&D~\S5.5.2),
$\mathrm{AICC} = \mathrm{AIC} + 2k(k+1)/(n - k - 1)$ with $k = p + q + 2$
counting the constant and $\sigma^2$. AICC values (params $= p+q$):
ARMA(2,1) $-8083.73$~$\checkmark$, ARMA(3,0) $-8083.68$,
ARMA(1,1) $-8081.53$~$\checkmark$, ARMA(2,0) $-8075.06$, ARMA(1,0) $-8036.14$.
ARMA(2,1) and AR(3) are essentially tied ($\Delta\mathrm{AICC} = 0.05$, well
below the rule-of-thumb threshold of 2). We pick \textbf{ARMA(2,1)} as the
main model and use AR(3) as a Yule--Walker check.

\textbf{Chosen model.} ARMA(2,1) MLE: $\hat\phi_1 = +0.441~(0.094)$,
$\hat\phi_2 = +0.279~(0.082)$, $\hat\theta_1 = +0.512~(0.091)$,
$\hat\sigma^2 = 1.24\!\times\!10^{-3}$. The constant is $-0.0004~(0.004)$ and
not significant ($p = 0.93$), which makes sense because $\hat Y_t$ is
mean-zero by OLS construction. The AR roots are $z \in \{1.26, -2.84\}$ and
the MA root is $-1.95$, all outside the unit disc.
\textbf{Yule--Walker} (B\&D~\S5.1.1) on the AR(3) candidate gives
$\hat\phi^{\mathrm{YW}} = (+0.948, -0.206, +0.075)$,
$\hat\sigma^2 = 1.25\!\times\!10^{-3}$, which matches the AR(3) MLE
$(+0.954, -0.207, +0.071)$ to three decimals. The two estimators target the
same population quantity, so they agree.
\textbf{Residual diagnostics.} Ljung--Box on the ARMA(2,1) residuals with the
B\&D~\S5.3 correction $\mathrm{df} = h - p - q$ gives $Q(20) = 43.7$,
$\mathrm{df}=17$, $p = 4\!\times\!10^{-4}$, so the model fails the residual
whiteness test. The remaining correlation sits at multiples of 7
(Figure~\ref{fig:diag}c). This means the harmonic weekly basis catches the
deterministic weekday \emph{shape} but not the shared noise between same
weekdays. A SARIMA term at lag 7 would handle this, but it is outside the
course toolkit, so we report the failure rather than add more parameters.

\section{Forecasting}
We hold out the last 30 days (2020-09-01 to 2020-09-30) and refit ARMA(2,1)
on the first $n-30$ residuals. The $h$-step-ahead forecast and its prediction
MSE $v_n(h)$ are computed from the ARMA recursion (B\&D~\S3.3.1;
\texttt{statsmodels.get\_forecast} runs the innovations algorithm internally).
The Gaussian 95\% PI $\hat Y_{n+h} \pm 1.96\sqrt{v_n(h)}$ is then transformed
back to MW by adding the extrapolated $\hat m_t + \hat s_t^{(\mathrm{year})}
+ \hat s_t^{(\mathrm{week})}$ and exponentiating. The MW interval is
asymmetric (log-normal), but since the transformation is monotone the 95\%
coverage is preserved. Exponentiating the conditional mean gives the
conditional \emph{median} of $X_{n+h}$ (the mean differs by a factor
$\exp(v_n(h)/2) \leq 1.0025$, which is negligible).

\textbf{Skill (held-out 30 days).}
$\mathrm{RMSE} = 187.6$\,MW, $\mathrm{MAE} = 145.1$\,MW,
$\mathrm{RMSE}/\mathrm{std}(\mathrm{train}) = 0.097$. The empirical 95\% PI
coverage is $30/30$, which is consistent with the nominal level (note
$0.95^{30} \approx 0.21$). Forecast and actual values agree to within about
$250$\,MW on every test day (Figure~\ref{fig:forecast}). After a few days
the residual forecast decays towards zero, so the medium-range forecast is
mostly the deterministic seasonal and trend part.

\section{Reflection}
\textbf{Difficulties.} The main problem was the leftover lag-7 correlation in
$\hat Y_t$, which a low-order ARMA cannot fully remove. SARIMA would fix it
but is outside the course, so we just report the Ljung--Box failure.
\textbf{Robustness.} Results are not very sensitive to the model choice:
ARMA(2,1) and AR(3) give point forecasts that differ by at most 4\,MW per day
and PI widths by less than 1\%. The biggest forecast errors come from
holidays and weather events that the seasonal model does not see.
\textbf{Alternatives.} A SARIMA term at lag 7, exogenous regressors (e.g.\
holiday flags, heating-degree-days), or a period of 365.25 instead of 365
would all be reasonable next steps.
\textbf{Software.} Python\,3 with \texttt{statsmodels}\,0.14
(\texttt{ARIMA}, \texttt{yule\_walker}, \texttt{acorr\_ljungbox},
\texttt{plot\_acf/pacf}), \texttt{numpy}, \texttt{pandas}, and
\texttt{matplotlib}. The pipeline can be reproduced with
\texttt{jupyter nbconvert --execute project\_part\_a.ipynb}.

\FloatBarrier
\clearpage
\appendix

\begin{figure}[H]
\centering
\includegraphics[width=0.72\linewidth]{figures/fig_decomposition.png}
\caption{Classical decomposition (B\&D~\S1.5) of daily SE\_3 load,
2015-01-01 to 2020-09-30. (a)~Raw daily mean load $X_t$, $n = 2100$.
(b)~Log-load $\log X_t$ (grey) overlaid with the fitted deterministic component
$\hat m_t + \hat s_t^{(\mathrm{year})} + \hat s_t^{(\mathrm{week})}$ (red);
seasonals are harmonic regressions ($d=365, k=2$ and $d=7, k=3$) and the trend
is a quadratic OLS fit. (c)~Cleaning residuals $\hat Y_t$ entering the ARMA stage.}
\label{fig:decomposition}
\end{figure}

\begin{figure}[H]
\centering
\includegraphics[width=0.78\linewidth]{figures/fig_diagnostics.png}
\caption{Diagnostic plots. (a)~Sample ACF of cleaning residuals $\hat Y_t$
(lags 0--50, blue bands at $\pm 1.96/\sqrt{n}$, B\&D~\S1.4.1). Slow decay
$\hat\rho(1) = 0.83$ and the lag-7 spike $\hat\rho(7) = 0.30$ motivate an
AR-dominated ARMA fit. (b)~Sample PACF of $\hat Y_t$, largest spike at lag~1.
(c)~ACF of ARMA(2,1) residuals; the surviving correlation is concentrated at
multiples of~7 ($Q(20) = 43.7$, df $=17$, $p = 4\!\times\!10^{-4}$).
(d)~ARMA(2,1) residual time series; level and amplitude stable.}
\label{fig:diag}
\end{figure}

\begin{figure}[H]
\centering
\includegraphics[width=0.88\linewidth]{figures/fig_forecast.png}
\caption{30-day forecast on the held-out window 2020-09-01 to 2020-09-30,
ARMA(2,1) refit on the first $n - 30$ cleaning residuals. Blue: last 90~days
of training. Black markers: held-out actual daily means.
Red: $h$-step ARMA forecast (B\&D~\S3.3.1) inverted to MW by adding the
extrapolated trend and seasonals and exponentiating; red band is the 95\% PI
(log-normal in MW). RMSE $= 187.6$\,MW, MAE $= 145.1$\,MW, empirical PI
coverage $30/30$.}
\label{fig:forecast}
\end{figure}

\end{document}
